\section{绪论}

近年来，无人机系统从单机作业向多机协同、集群智能方向快速发展，已成为航空航天、人工智能、机器人工程等领域的研究热点。

\subsection{研究背景与意义}

伴随行业应用持续深入，单个无人机于作业范围、执行效率、任务冗余性等方面的局限越发显著，而无人机集群凭借并行作业、分布式感知、协同决策、高鲁棒性等优点，在民用与军用领域都展现出无法替代的价值。近些年来，无人机技术加速更新，单个无人机已然无法满足复杂场景的要求，无人机集群技术成为行业的研究重点。在军事侦察、边境巡查、灾害救援、物流运送、智慧城市测绘、电力巡查、环境监测等领域，无人机集群依靠并行作业、协同感知、冗余容错、高效覆盖等优势，逐渐取代传统单机作业模式，呈现出规模化、智能化、协同化的发展趋势\cite{wu2003}。无人机集群具有重要价值，它借助多机之间的信息交互、任务分配、轨迹协同、姿态配合等方式，达成单机难以实现的大规模、高精度且高鲁棒性的作业任务。

当下主流的无人机仿真平台有ROS/Gazebo、PX4 SITL、MATLAB/Simulink、Webots、微软推出的AirSim。当中，AirSim是基于Unreal Engine或者Unity引擎来进行开发的，它具有开源免费的特点，物理仿真精度比较高，视觉渲染效果非常逼真，能够支持多旋翼、固定翼、无人车等多种载体，并且能很好地兼容ROS机器人操作系统\cite{airsim}，其提供了Python/C++双接口，被全球范围内的高校、科研机构、企业广泛应用于无人机感知、控制、规划、集群算法验证等方面。

不过在实际的工程应用里，AirSim存在着一些性能方面的缺陷：当仿真无人机数量比较少（1至10架）的时候，系统运行比较流畅，帧率相对稳定，时延也比较低。当数量增长到20架以上时系统性能发生急剧下滑，具体体现为帧率从60FPS以上快速下降到10FPS以下，仿真时延从毫秒级别提高至百毫秒级别，CPU占用率一直处于满负荷状态，内核态开销出现异常增高，线程频繁进行切换，仿真出现失步情况，控制指令产生延迟，甚至会出现引擎崩溃等状况。这一问题对AirSim在大规模无人机集群算法验证方面的应用造成了严重限制，使得研究者不得不采用小规模仿真或者简化模型，无法如实反映集群系统在大数量级时的协同特性、通信冲突、调度瓶颈、资源竞争等实际行为，大幅降低了仿真结果的可信度、工程参考价值\cite{chen2022,wu2023,dong2024}。

从技术的本质角度来分析，AirSim所存在的性能瓶颈并非是由于硬件算力方面有不足而导致的，而是软件方面的线程调度机制不够合理，同步等待逻辑效率较低，并且内核态与用户态切换过于频繁、自旋等待策略在超线程CPU环境下失效等底层问题。目前国内外相关研究多集中于利用AirSim进行算法验证\cite{fang2025}，极少针对其底层并发机制、线程模型、定时逻辑开展系统性性能优化\cite{lin2022,han2023,gao2022}，导致大规模集群仿真的性能难题长期未得到解决。

在这种背景状况下，本课题将基于AirSim的无人机集群仿真软件性能优化当作研究目的，重点关注线程调度、内核开销、等待机制这三个核心问题，借助理论剖析、瓶颈查找、方案规划、源码修改、实验验证等一系列操作，建立一套完备的性能优化模式，大幅度提高AirSim在大规模集群场景中的运行效率、稳定性\cite{chen2023,sun2023,hu2022}，为无人机集群算法的研发给予了高效且可靠的仿真支持。

\begin{figure}[H]
    \centering
    \includegraphics[width=0.9\textwidth]{figures/image2.png}
    \caption{无人机集群典型应用场景图}
    \label{fig:scene}
\end{figure}

本研究围绕AirSim大规模无人机集群仿真性能瓶颈展开深度优化\cite{zhu2021}，兼具理论意义、应用意义与实践意义，对推动无人机集群技术发展、完善仿真平台性能、提升工程研发效率具有重要价值。

\textbf{（1）理论意义}

第一，本研究深入揭示高并发实时仿真系统的性能退化机理，将操作系统线程调度理论、内核态/用户态切换开销理论、自旋等待/阻塞睡眠理论与无人机集群仿真工程实践结合，建立多实体并发仿真的性能分析模型，完善高性能仿真领域的理论体系。

第二，研究明确超线程CPU环境下，线程数量超过物理核心数时，低效sleep(0)、yield()等调用引发"调度风暴"的内在规律，提出线程等待模型的优化准则，为高并发实时系统性能建模提供新视角。

第三，构建"问题复现---瓶颈定位---方案设计---实验验证---迭代优化"的标准化仿真平台调优理论框架\cite{wu2023}，为数字孪生、机器人仿真、实时渲染、自动驾驶仿真等领域提供可复用的理论范式。

\textbf{（2）应用意义}

第一，优化后的AirSim可稳定支持30架及以上无人机集群仿真，帧率、时延、稳定性大幅提升，能够满足集群控制、路径规划、协同感知、避障、通信组网等算法的大规模验证需求，提高仿真结果真实性与可信度。

第二，降低无人机集群算法研发成本，减少真机试飞次数，规避安全风险，缩短算法迭代周期，加速无人机集群技术在物流、救援、测绘、巡检、军事等领域的工程落地。

第三，优化方案以源码改造形式实现，不依赖硬件升级，兼容现有Windows/Linux系统与Unreal/Unity引擎，具备良好的通用性与可移植性，可直接被全球研究者使用。

\textbf{（3）实践意义}

第一，形成一套面向AirSim的系统化性能调优流程，包括环境搭建、问题复现、瓶颈定位、工具使用、源码修改、测试评估等完整工程实践方法，可为同类仿真平台优化提供直接参考。

第二，通过对AirSim核心C++源码的深度解析与改造\cite{li2023}，提升机器人工程、计算机科学等专业学生的大型工程源码阅读、并发编程、系统调优、性能测试等工程实践能力。

第三，研究成果可集成到开源AirSim项目，贡献社区，提升国产优化方案在国际开源领域的影响力，推动无人机仿真技术的共同进步。

\subsection{国内外研究现状}

国外无人机集群仿真研究起步较早，以微软、NVIDIA、Intel等企业主导的仿真平台为核心，高校与科研机构聚焦算法验证。

在性能优化相关领域里，国外学者着重研究实时系统线程调度、同步机制。Lamport提出了经典线程同步理论，这为锁机制、条件变量还有信号量奠定了基础，Butenhof深入剖析了POSIX线程编程，指出了自旋等待与阻塞睡眠的适用边界所在，Intel研究人员指出，在超线程CPU当中，当线程数超过物理核心数的时候，yield()与sleep(0)会致使上下文切换开销急剧增加，内核态CPU占比出现异常升高的情况。

在集群仿真方面，斯坦福大学、麻省理工学院、加州大学伯克利分校等机构借助AirSim展开无人机控制、避障、视觉SLAM、强化学习研究，不过仿真规模一般控制在10架以内，并没有触及大规模性能瓶颈。部分学者采用分布式仿真，把多机负载分散到多台计算机上，然而该方案部署繁杂、通信延迟大、成本高，无法替代单机高性能仿真。国外研究拥有技术先发优势，在仿真平台建立、基础并发理论方面取得了丰富成果，但是缺少针对AirSim大规模集群场景的底层线程优化，没有形成系统化、能够落地的性能优化方案。

国内无人机集群仿真研究近年来发展迅速，聚焦平台构建、算法验证、工程应用，性能优化研究逐步受到关注。

在仿真平台方面，冯聪（2018）设计交互式无人机集群编队仿真平台，强化人机交互与可视化\cite{feng2018}；方仪豪、邹丹平（2025）在多旋翼无人机仿真平台综述中指出，AirSim在视觉保真度、API灵活性方面优势显著，但多机并发效率不足，是制约大规模仿真的关键因素。

在AirSim应用方面，郑筱宇等（2025）基于AirSim实现无人机位置跟踪与编队控制；方璨琦（2023）利用AirSim开展深度强化学习路径规划；周游（2020）基于AirSim研究三维避障算法；刘彦辰等（2024）研究分布式集群避碰规划；石辅天等（2023）、朱博顺等（2021）分别开展多无人机协同搜索、蜂群侦察仿真研究。上述研究均验证了AirSim的有效性，但仿真规模普遍小于15架，未解决大规模性能下降问题。

在性能优化方面，冯一飞（2023）采用分布式架构分摊负载，提升仿真规模；郭首江（2021）从网络协议层面优化集群接入性能；龙腾飞（2023）研究AirSim虚实交互平台实时性改进；章铭泽（2025）关注集群组网通信仿真效率。这些研究从不同角度提升仿真性能，但均未针对AirSim内部线程等待、调度机制、内核态开销等核心瓶颈开展深度优化。

国内研究充分认可AirSim的价值，也明确大规模集群仿真存在性能痛点。不过，根因定位不够清晰，优化方案不够系统，且缺乏多方案对比验证，这与工程化、标准化、可推广的优化目标存在较大差距。

本研究具有鲜明创新性与必要性，以操作系统并发理论为指导，将AirSim源码作为改造对象，把大规模集群仿真设为目标，去定位性能瓶颈，设计多套优化模型，进行量化测试，优选最优方案，填补国内外研究空白，为无人机集群仿真提供高效可靠的技术支撑。

\subsection{研究内容与技术路线}

本研究聚焦于AirSim无人机集群仿真性能的优化工作，依照"理论分析→瓶颈定位→方案设计→实现改造→实验验证→总结归纳"这样的技术路径来推进，具体的研究内容呈现如下：

首先是理论基础与文献研究部分，对无人机集群技术、仿真平台分类、AirSim架构、线程调度、同步等待、内核态开销、性能监测工具等一系列相关理论和文献进行系统梳理，以此明确研究的边界范围、性能瓶颈的指向方向、优化所依据的理论基础。

其次是仿真环境搭建与性能问题复现环节，建立起统一的硬件与软件测试环境，借助Unity搭建出能够进行扩展的多无人机集群仿真场景，编写Python自动化控制脚本，重现当20架以上无人机进行仿真时出现帧率急剧下降、时延明显增高、CPU异常占用等性能问题的情况，从而获取基础性能数据，并确定性能的拐点位置。

最后是性能瓶颈深度定位与根因分析阶段，运用Windows Performance Analyzer、Linux perf、自定义监测模块，从CPU占用、线程状态、内核态与用户态占比、热点函数、调用栈、上下文切换频率等多个维度入手，找出AirSim在大规模集群环境下的性能瓶颈所在，明确低效线程等待、调度风暴、超线程冲突乃是其核心的根因所在。

\subsection{本章小结}

本章从无人机集群技术发展与仿真平台需求出发，阐述研究背景与意义；系统梳理国内外研究现状，明确当前研究不足与本研究创新价值；确定研究内容、技术路线与论文组织结构，为后续研究提供清晰框架与方向。
